\chapter{Үр дүн ба үнэлгээ}

\section{Туршилтын орчин}
Бид ML загварыг сургахдаа болон FastAPI серверийг гүйцэтгэхдээ 8GB-аас дээш санах ойтой, зохих түвшний график боловсруулах процессор (GPU) бүхий орчинд ажиллаж тестчилсэн. Серверийн туршилтыг Docker хэрэгслээр контейнержуулж, MongoDB-тэй холбогдох асуудлыг урьдчилан шийдвэрлэсний үндсэн дээр тогтвортой ажиллах орчныг бүрдүүлсэн.

\section{Загварын нарийвчлалын хэмжүүрүүд}
Объект илрүүлэлтийн загварын үнэлгээнд түгээмэл ашиглагддаг нийтлэг параметр болох Дундаж Нарийвчлал (Mean Average Precision - mAP), Нарийвчлал (Precision), Мэдрэг байдал (Recall) зэргийг тус тус хэмжсэн.
\begin{itemize}
    \item \textbf{Precision (Нарийвчлал):} Нийт эерэг гэж таамагласан объектуудаас хэд нь үнэхээр тухайн ангиллын бараа байсныг илтгэнэ.
    \item \textbf{Recall (Мэдрэг байдал):} Бодит зураг дээр байгаа нийт объектуудын хэдэн хувийг нь зөв илрүүлж чадсаныг заана.
    \item \textbf{mAP50:} Итгэлцлийн интервал (IoU threshold) 0.5 үеийн нарийвчлалын дундаж.
\end{itemize}

Туршилтын үр дүнд манай сургасан YOLOv8n загварын хувьд Validation багц дээр mAP50 нарийвчлал нь маш өндөр ($>90\%$) гарсан. Энэ нь загвар лонхны хэлбэр, мөн хаяг шошгыг сайн сурсан бөгөөд барааг өндөр нарийвчлалтайгаар таньж байна гэсэн үг юм. Загвар нь давхарлаж өрсөн, эсвэл гэрэлтүүлэг хангалтгүй нөхцөлд ч алдаа багатай илрүүлэлт хийсэн.

\section{Системийн хурдны үнэлгээ ба Бодит ажиллагаа (Real-time constraints)}
Зөвхөн нарийвчлал сайн байхаас гадна, боловсруулалт хурдан байх нь системд нэвтрүүлэх том шалгуур байдаг. Сервер талд YOLOv8n хөнгөн загвар ашигласны давуу тал гарч, нэг зургийг илрүүлэхэд ойролцоогоор $200-300$ миллисекунд хугацаа зарцуулагдаж байгаа статистик гарсан. 
Иймд энэхүү хугацаа нь хүлээгдэх хугацааны шаардлагыг бүрэн хангаж байгаа ба сүлжээний саатал (latency)-ийг оролцуулаад нийтдээ 1 секунд орчмоор хэрэглэгчид аудитын хариу очиж буй нь их сургуулийн судалгааны болон салбарын бодит шаардлагуудад бүрэн нийцсэн үзүүлэлт болж чадсан \cite{yolov8_ultralytics}.
